\chapter{Beneficios que aporta el trabajo y relación con los Objetivos de Desarrollo Sostenible}
\label{ch:beneficios}

\section{Beneficios del sistema desarrollado}

El sistema desarrollado convierte un archivo de casi veinte años de actas
municipales --en la práctica, solo accesible leyendo PDF uno por uno-- en una
fuente consultable en lenguaje natural, con respuestas respaldadas por cifras
verificables y citas a la página exacta del documento de origen. Esto beneficia,
al menos, a tres perfiles distintos:

\begin{itemize}
	\item \textbf{Ciudadanía}: permite conocer de forma directa qué se ha
	debatido y votado sobre un tema o un grupo político, sin depender de
	resúmenes de terceros ni de la capacidad de leer cientos de páginas de actas.
	\item \textbf{Periodistas e investigadores}: permite hacer preguntas
	agregadas --evolución de un tema por año, comparación entre grupos-- que de
	otro modo exigirían un trabajo manual de cómputo sobre el archivo completo.
	\item \textbf{La propia administración municipal}: demuestra, con datos
	reales de veinte años de plenos, que la información pública ya generada por
	el propio Ayuntamiento puede reutilizarse para mejorar su accesibilidad sin
	coste de recogida de datos adicional.
\end{itemize}

Adicionalmente, la comparación sistemática entre RAG vectorial y GraphRAG sobre
un mismo corpus real --y no sobre un \textit{benchmark} sintético-- aporta
evidencia empírica sobre en qué tipo de preguntas destaca cada arquitectura,
resultado con valor más allá de este proyecto concreto.

\section{Relación con los Objetivos de Desarrollo Sostenible}

Este trabajo se relaciona principalmente con dos Objetivos de Desarrollo
Sostenible (ODS) de la Agenda 2030 de Naciones Unidas:

\begin{itemize}
	\item \textbf{ODS~16 (Paz, justicia e instituciones sólidas)}, en particular
	la meta 16.10, que persigue garantizar el acceso público a la información. El
	sistema reduce la barrera práctica de acceso a la actividad de un órgano de
	gobierno local, facilitando la transparencia y la rendición de cuentas ante
	la ciudadanía.
	\item \textbf{ODS~11 (Ciudades y comunidades sostenibles)}, en la medida en
	que contribuye a una gestión municipal más accesible y participativa,
	facilitando que cualquier persona pueda informarse sobre las decisiones que
	afectan a su ciudad.
\end{itemize}
